% Full JACT-style manuscript
\documentclass[11pt]{article}
\usepackage{amsmath,amssymb,amsthm}
\usepackage{hyperref}
\usepackage{graphicx}
\usepackage{geometry}
\usepackage{booktabs}
\usepackage{fancyvrb}
\geometry{margin=1in}

\title{Obstruction Persistence and Witness Failure in Finite Regional Cohomology}
\author{Duston Moore\\Independent Researcher, Edmonton, Canada\\dhwcmoore@gmail.com}
\date{June 2026}

\begin{document}
\maketitle

\begin{abstract}
We present a finite cohomological obstruction classifier with machine‑readable
certificates for seam residues in finite regional covers. The framework
distinguishes direct obstruction persistence in refined nerves from witness
 (cycle) persistence, and characterizes cycle‑faithfulness via rational rank
tests. We give algorithms to produce audit‑grade JSON certificates and
rank/failure witnesses, prove correctness of the tests, and demonstrate
persistence (or its certified failure) across several refinement strategies.
All code, inputs, and certificates are provided for full reproducibility.
\end{abstract}

\section{Introduction}
We study finite regional cohomological obstructions arising from seam residues
in combinatorial covers. The main contributions are:
\begin{itemize}
  \item A practical classifier that tests whether a residue is a cocycle and
    whether it is a coboundary; the implementation emits machine‑readable JSON
    certificates that include cohomological summaries and cycle witnesses.
  \item A distinction between two persistence notions under refinement:
    direct persistence (a refined cycle pairs nontrivially with the transferred
    residue) and witness persistence (the original witness cycle lifts to the
    refined nerve up to a nonzero rational factor). We prove algebraic
    criteria distinguishing them.
  \item Rational rank tests (nullspace and rank criteria) that certify
    cycle‑faithfulness or provide failure witnesses via non‑uniform edge ratios.
\end{itemize}

This paper is accompanied by a reproducible repository containing the code,
the JSON inputs and certificates, and scripts to reproduce all experiments.

\section{Background and notation}
Let $N$ be a finite oriented nerve with regions $U_i$. We work over $\mathbb{Q}$.
Denote cochain groups by $C^k(N;\mathbb{Q})$ with coboundary operators
$\delta^k:C^k\to C^{k+1}$. For a residue $r\in C^1$, we test the cocycle
condition $\delta^1 r = 0$ and whether $r$ lies in the image of $\delta^0$.
The cohomology group is $H^1 = \ker\delta^1/\operatorname{im}\delta^0$.

\section{Classification algorithm}
The classifier proceeds in two steps:
\begin{enumerate}
  \item Cocycle test: verify $\delta^1 r = 0$ (trivial when $C^2=0$).
  \item Coboundary test: solve the rational linear system $\delta^0 b = r$.
    If no solution exists, the residue represents a nonzero $H^1$ class.
\end{enumerate}

Implementation details: we use exact rational arithmetic (SymPy) to avoid
floating point error. The classifier emits a JSON certificate with fields
including `cohomology_summary`, `cycle_witness` (if available), and an audit trail.

\section{Invariance under presentation}
We prove the classification is invariant under the five tested presentation
changes: region renaming, edge orientation reversal, nonzero rational scaling,
gauge perturbation ($r\mapsto r+\delta^0 b$), and edge order permutation.
Each operation induces an isomorphism or leaves cohomology classes unchanged.
The implementation verifies these operations on the base object and reports
`stable_under_tested_presentations` in the invariance report.

\section{Refinements and transfer maps}
We consider refinements $\rho: N'\to N$ with transfer maps $\rho^*:C^1(N)\to C^1(N')$.
Our primary transfer strategy is equal‑distribution: split an old edge value
equally across the refined edges replacing it; internal refined edges receive 0.

Two persistence notions:
\begin{description}
\item[Direct persistence:] There exists $z'\in Z_1(N')$ with $\langle z',\rho^* r\rangle\ne0$.
\item[Witness persistence (cycle‑faithfulness):] There exists $z'\in Z_1(N')$ and
  $\lambda\in\mathbb{Q}^\times$ with $\rho_* z' = \lambda z$ for an original
  cycle $z$ pairing nontrivially with $r$.
\end{description}

We prove the Cycle‑Lift Persistence Theorem: if $\rho_* z' = \lambda z$ then
$[\rho^* r]\ne0$ whenever $\langle z,r\rangle\ne0$, by adjointness of $\rho^*$ and $\rho_*$.

\section{Cycle‑faithfulness and rational rank tests}
Cycle‑faithfulness can be decided by two algebraic methods:
\begin{itemize}
\item Nullspace method: solve a linear system for $(z',\lambda)$ certifying a nonzero lift.
\item Rank criterion: let $K$ be a basis for $Z_1(N')$ and $P=\rho_*$. Then $\rho$ is
  cycle‑faithful iff $\operatorname{rank}(PK)=\operatorname{rank}([PK\; z])$.
\end{itemize}

When $\operatorname{rank}(PK)=1$, faithfulness reduces to checking uniformity
of ratios $v_e/z_e$ for a generator $v$ of $\operatorname{im}(PK)$.
Failure is certified by a non‑uniform ratio vector (a short, exact rational witness).

\section{Experiments and certificates}
We applied the classifier to the canonical examples and to the actual object
`actual_gluing_object_v1`. The base classification and certificate summary:

\begin{table}[h]
\centering
\begin{tabular}{l l}
\toprule
Item & Value \\
\midrule
Classification & nontrivial\_H1\_obstruction \\
Is cocycle & true \\
Is coboundary & false \\
Pairing (cycle witness) & -5 \\
$\dim H^1$ & 1 \\
\bottomrule
\end{tabular}
\caption{Base object classification summary (from certificate).}
\end{table}

Refinement experiment summary (pairings shown):
\begin{table}[h]
\centering
\begin{tabular}{l r r}
\toprule
Refinement & Pairing & Obstruction persists? \\
\midrule
Subdivide $U_1$ & -7/2 & Yes \\
Subdivide $U_2$ & -4 & Yes \\
Subdivide all regions & -5/4 & Yes \\
Insert bridge $U_1$–$U_2$ & -5 & Yes \\
\bottomrule
\end{tabular}
\caption{Refinement pairings (direct persistence).}
\end{table}

Implementation and reproducibility: the repository contains the code used to
generate the certificates and reports. Full JSON certificates are archived in
`certificates/`; a summary of their key fields appears in Appendix~\ref{app:certs}.

\section{Discussion}
The experiments highlight an important distinction: obstruction persistence is
more robust than witness persistence. Equal‑distribution transfer maps can
preserve a nonzero cohomology class while making the original cycle fail to lift
proportionally. The rank criterion and failure witnesses provide exact, checkable
certificates for both outcomes.

\section*{Acknowledgements}
This work started as an independent project and benefited from feedback during
informal seminars. The author thanks colleagues for suggestions on reproducibility.

\appendix
\section{Summary tables of certificates}
\label{app:certs}
Key certificate fields (base object):
\begin{tabular}{ll}
Residue & (1,1,1,-2) \\
Support & U1-U2, U2-U3, U3-U4, U1-U4 \\
Pairing & -5 \\
dim C0, C1 & 4, 4 \\
dim H1 & 1 \\
\end{tabular}

\section{Full proofs (from repository)}
\label{app:proof}
The complete formal proofs are reproduced from `PROOF.md` for completeness.
\VerbatimInput{../PROOF.md}

\end{document}
